\section{Starting point: GRAILS}
\label{sec:grails}

This section summarises GRAILS~\citep{kulkarni2025grails} and the argument it
makes, because the rest of the paper follows the same argument into a setting
with several data sources.

\subsection{The problem GRAILS states}

GRAILS observes that ethical compliance in software is usually reached by
people who follow guidelines, regulations and other documents by hand, and
that there are few software architectures or tools that automate compliance
inside an application. It describes two existing routes and the cost of each:
\begin{itemize}
  \item \textbf{Privacy by Design}~\citep{cavoukian2010privacy} and related
        design methods, such as Value Sensitive Design, embed privacy early in
        the development life cycle. GRAILS finds them suited to new systems and
        notes that they increase design complexity.
  \item \textbf{Retrofitting} adds privacy to an existing system, by
        redevelopment, by modification, or by \emph{co-existence}: adding
        measures without changing the core functions. GRAILS notes that the
        first two need large changes and that retrofits can add technical debt
        and performance overhead.
\end{itemize}
It also argues that privacy impact assessments, data loss prevention systems
and access-control models such as role-based access control mostly reduce
privacy risk after the software is built, that model-view-controller designs
and the usual development life cycle have no place for ethical constraints,
and that specialised privacy tools are often specific to one domain.

\subsection{The framework}

GRAILS has two parts (\cref{fig:grails}):
\begin{itemize}
  \item The \textbf{Ethical-Restriction Rails} decide. A user request is a tuple of the user, the data request and its
        context.
        A \emph{Know Your User} (KYU) service turns user attributes (user type,
        operation, information classification, data ownership and compliance
        history) into a trust score, using fuzzy inference in the proof of
        concept. A \emph{Data Protection Service} builds a data protection
        profile from file-level and item-level metadata and from a knowledge
        base of regulations; the proof of concept uses a large language model
        fine-tuned on Indian regulatory frameworks and healthcare data. A
        \emph{Restriction Planner}, which GRAILS compares to a database query
        processor, maps the KYU profile and the data protection profile to
        restriction strategies taken from a \emph{Strategy Knowledge Base}.
  \item The \textbf{Ethical Guard} enforces. Its \emph{Strategy Engine} turns
        the restriction plan into an execution plan, and its \emph{Execution
        Engine} applies the plan to data held in a \emph{Privacy Vault}, through
        built-in implementations, external services or plug-ins, and returns
        the filtered data.
\end{itemize}

Formally, GRAILS writes the filtered output as
\begin{equation}
  D_{\text{out}} = \bigcup_{g \in G} T_g(D_g, \lambda, \nu),
  \qquad G = \{\text{cell}, \text{row}, \text{column}, \text{table}\},
\end{equation}
where \(D_g\) is the projection of the dataset at granularity \(g\),
\(\lambda \in \{\text{low}, \text{moderate}, \text{high}\}\) is the privacy level
and \(\nu \in \{\text{low}, \text{moderate}, \text{high}\}\) is the user's trust
level. Its Table~1 gives the strategies for each of the 36 combinations of
granularity, sensitivity and trust, drawn from statistical disclosure control:
cell suppression, top and bottom coding, microaggregation, generalisation,
binning, full and partial masking, noise injection and differential privacy.

GRAILS measures how much a response was filtered with a \emph{Filtering Score},
the mean over all cells of a penalty \(\delta \in [0,1]\) between the original
and the filtered value. It evaluates the framework on the \emph{Karnataka At a
Glance} open government dataset, which covers 30 districts and 17 domains with
more than 1,200 attributes. It reports that filtering is highest where
sensitivity is high and trust is low, and lowest where sensitivity is low and
trust is high. The proof of concept is written in Java and uses the Facade,
Decorator and Strategy design patterns over tabular data in an SQL database.

\begin{figure}[t]
\centering
\begin{tikzpicture}[
  node distance=4mm and 6mm,
  box/.style={draw=dgInk, rounded corners=2pt, align=center, font=\footnotesize, inner sep=4pt},
  part/.style={draw=dgRule, dashed, rounded corners=3pt, inner sep=5pt},
  arr/.style={-{Stealth[length=2mm]}, draw=dgInk},
  lbl/.style={font=\scriptsize\itshape, text=dgInk}
]
\node[box] (app) {Application\\software};
\node[box, right=10mm of app] (exec) {Execution\\Engine};
\node[box, below=of exec] (strat) {Strategy\\Engine};
\node[box, right=of exec] (vault) {Privacy\\Vault};
\node[box, right=14mm of vault] (kyu) {KYU\\service};
\node[box, right=of kyu] (dps) {Data Protection\\service};
\node[box, below=8mm of $(kyu)!0.5!(dps)$] (rp) {Restriction\\Planner};
\node[part, fit=(exec)(strat)(vault), label={[lbl]above:Ethical Guard (runtime)}] (guard) {};
\node[part, fit=(kyu)(dps)(rp), label={[lbl]above:Ethical-Restriction Rails}] (rails) {};
\draw[arr] (app) -- node[lbl, above] {request} (exec);
\draw[arr] (exec.east) -- (vault.west);
\draw[arr] (guard.east |- kyu) -- node[lbl, above] {request} (kyu.west);
\draw[arr] (kyu) -- (rp);
\draw[arr] (dps) -- (rp);
\draw[arr] (rp.west) -- node[lbl, below] {restriction plan} (strat.east);
\draw[arr] (strat) -- node[lbl, left] {execution plan} (exec);
\end{tikzpicture}
\caption{The two parts of GRAILS, simplified from Figure~1 of
\citet{kulkarni2025grails}. The Rails decide; the Guard enforces, on data held
in one Privacy Vault.}
\label{fig:grails}
\end{figure}

\subsection{What GRAILS leaves open}
\label{sec:grails-open}

GRAILS states its own next steps: other ethical dimensions such as fairness and
transparency, learned models for the KYU profile and the data protection level,
and retrieval over ontologies of regulations. Reading the paper for the setting
of this survey, where data sits at several sources and models are trained
across them, we see five further open points. They are not defects of GRAILS,
which addresses one application and one store.
\begin{enumerate}[label=O\arabic*.]
  \item \textbf{One boundary.} The Guard filters the response to a request. In
        a federated study, data also leaves a site as a model update, a metric
        or a count, and a model trained on it may be released. These are
        further boundaries.
  \item \textbf{One store.} The Guard reads a central Privacy Vault. When each
        source must keep its records, the Guard has to run at each source.
  \item \textbf{Reproducible decisions.} The proof of concept derives the data
        protection profile with a fine-tuned language model. A decision that a
        learned model makes may differ between two runs or two model versions,
        which makes it harder to replay during an audit.
  \item \textbf{Failure behaviour.} The paper does not state what the Guard does
        when the Rails cannot answer.
  \item \textbf{Utility.} The Filtering Score measures how much was changed. It
        does not measure how much an analysis or a model loses because of the
        change.
\end{enumerate}
\Cref{sec:design} describes how Dagents addresses O1 to O4, and
\cref{sec:case-study} measures O5 on one study.
